\documentclass[11pt]{article}
\usepackage[margin=1in]{geometry}
\usepackage{amsmath,amssymb,booktabs,array}
\usepackage[T1]{fontenc}
\usepackage{lmodern}
\usepackage{xcolor}
\usepackage{hyperref}
\hypersetup{colorlinks,linkcolor=black,urlcolor=blue!60!black}
\setlength{\parskip}{4pt}
\title{ShiftAddQwen mat-vec: symbols and the sign-select identity}
\author{ShiftAddQwen notes}
\date{October 2026}
\begin{document}
\maketitle

\section{Setting}
A linear layer computes $o = W x$. The ShiftAdd format (ggml types \texttt{sa\{2,\dots,6\}[\_128|\_256]}) stores each
row of $W$ as a sum of \emph{sign planes} with one scale per plane and per group of weights along the input dimension:
\begin{equation}
w_j \;\approx\; \sum_{p=1}^{Q} \alpha_p \, b_{p,j}, \qquad b_{p,j} \in \{-1,+1\}.
\label{eq:bcq}
\end{equation}
The decode mat-vec in the patched \texttt{llama.cpp} evaluates one output row against one input block of this form. There
are no lookup tables; the weights only ever appear as bit masks.

\section{Symbols}
\begin{center}
\renewcommand{\arraystretch}{1.25}
\begin{tabular}{>{$}l<{$} p{0.78\textwidth}}
\toprule
\text{symbol} & meaning \\
\midrule
j = 1,\dots,G & index of a weight inside one \emph{group} along the input dimension; $G \in \{64, 128, 256\}$ weights share one scale per plane \\
w_j & the original (bf16) weight at position $j$ of the output row \\
p = 1,\dots,Q & index of the sign plane; $Q$ planes = $Q$ bits per weight ($Q = 2$ to $6$ in the released models) \\
b_{p,j} \in \{-1,+1\} & the sign of weight $j$ in plane $p$ \\
s_j & shorthand for $b_{p,j}$ when a single plane $p$ is fixed (as in the kernel comment) \\
\mathrm{bit}_{p,j} \in \{0,1\} & how the sign is stored: $\mathrm{bit}=1 \Leftrightarrow s=+1$, $\mathrm{bit}=0 \Leftrightarrow s=-1$, i.e.\ $s_j = 2\,\mathrm{bit}_{p,j} - 1$ \\
\alpha_p & scale of plane $p$ for this group: $\alpha_p = d \cdot a_p$ with one fp16 $d$ per 256-weight super-block and one u8 code $a_p$ per plane and group \\
x_j & the input activation at position $j$ (fp16 in the model) \\
y_j & the activation after llama.cpp's q8 activation quantization: an int8 value, with one scale $\delta$ per block of 32 activations, $x_j \approx \delta\, y_j$ \\
\bottomrule
\end{tabular}
\end{center}

\section{The identity}
The contribution of plane $p$ of one group to the output row is $\alpha_p \sum_j s_j\, y_j$. Substituting
$s_j = 2\,\mathrm{bit}_j - 1$:
\begin{equation}
\sum_{j=1}^{G} s_j\, y_j \;=\; \sum_{j=1}^{G} (2\,\mathrm{bit}_j - 1)\, y_j
\;=\; 2 \sum_{j:\,\mathrm{bit}_j = 1} y_j \;-\; \sum_{j=1}^{G} y_j .
\label{eq:identity}
\end{equation}
Both terms are evaluated with \texttt{dp4a}, a single CUDA instruction (available since compute capability 6.1, so on
the V100 and A100): given two 32-bit registers $a$ and $b$ that each pack four int8 values, and a 32-bit accumulator
$c$, it computes
\begin{equation}
\texttt{dp4a}(a, b, c) \;=\; c + a_0 b_0 + a_1 b_1 + a_2 b_2 + a_3 b_3
\label{eq:dp4a}
\end{equation}
in one cycle on the integer units. llama.cpp's own k-quant mat-vecs are built on it. In the ShiftAdd kernel $b$ holds
four activations $y_j$ and $a$ holds four bytes that are each 0 or 1 (one nibble of sign bits expanded to bytes), so
the ``multiply'' is a select: a 1 byte passes the activation, a 0 byte drops it. With $a = \texttt{0x01010101}$ it is
a plain four-way sum. The two terms of Eq.~\eqref{eq:identity} then have different costs:
\begin{itemize}
\item $\sum_j y_j$ does not depend on the plane or on the output row. It is computed once per activation block and reused
  for every plane and every row (\texttt{sx} in the kernel, one \texttt{dp4a} against \texttt{0x01010101}).
\item $\sum_{\mathrm{bit}_j=1} y_j$ is a sum of activations \emph{selected} by the bit mask. Four sign bits (one nibble)
  are expanded into four 0/1 bytes and one \texttt{dp4a} against the four int8 activations accumulates exactly the
  masked sum. No multiplication by a weight value ever happens.
\end{itemize}

\section{Full output}
For one output row, summing over groups $g$ and planes $p$, and restoring the activation scale $\delta$,
\begin{equation}
o \;\approx\; \sum_{g} \sum_{p=1}^{Q} \alpha_{g,p}\,
\Big( 2 \sum_{j \in g:\, \mathrm{bit}_{g,p,j}=1} y_j \;-\; \sum_{j \in g} y_j \Big) \cdot \delta .
\end{equation}
The scale $\alpha_{g,p}$ is applied once per group and plane (an fp32 multiply-add on the integer partial sum), so the
inner loop is purely integer: per plane and per four weights, one \texttt{dp4a} plus a few integer ops to form the byte
mask. Prefill (batch $>$ 1) takes a different path: the planes are dequantized to fp16 with Eq.~\eqref{eq:bcq} and the
GEMM goes through cuBLAS.

\section{Relation to ShiftAddLLM}
ShiftAddQwen keeps the weight representation and the goal of ShiftAddLLM (NeurIPS 2024) and replaces the fitting
pipeline and the kernels.

\paragraph{Kept.}
\begin{itemize}
\item The binary-coding reparameterization of Eq.~\eqref{eq:bcq}: every weight is a sum of signed planes with
  per-group scales, so a mat-vec needs no weight multiplications.
\item The block-wise and column-wise scale layouts of the paper (the \texttt{sab<Q>\_<bc>x<br>} ggml family covers both),
  and mixed bit allocation across layers.
\end{itemize}

\paragraph{Replaced.}
\begin{itemize}
\item \textbf{Kernel.} ShiftAddLLM's kernels precompute, for every $k$ consecutive activations, the $2^k$ possible signed
  partial sums in a lookup table (LUT), index it with each plane's bit pattern, and handle the scales with shifts.
  ShiftAddQwen's decode kernel uses Eq.~\eqref{eq:identity} instead: the shared term $\sum_j y_j$ removes the sign
  handling and a dp4a byte mask does the selection, so there is no LUT storage, no LUT build, and the scale is one
  fp32 multiply-add per group and plane. The arithmetic that remains is select, integer add and a per-group scale.
  Prefill dequantizes to fp16 and uses cuBLAS.
\item \textbf{Fit.} ShiftAddLLM fits the planes in weight space with its reparameterization-aware objective.
  ShiftAddQwen fits each linear to the layer's output (GPTQ column sweep with act-order, exact Hessian-weighted scale
  refit), chooses the bits and group size per linear from measured end-to-end KLD sensitivities, and then distils every
  scale against the BF16 teacher (Section~6.2). The last two steps have no counterpart in the paper.
\item \textbf{Target and integration.} The paper evaluated dense OPT / LLaMA models in PyTorch with custom CUDA.
  ShiftAddQwen targets Qwen3.5's hybrid Gated-DeltaNet / attention architecture with its MTP head, keeps the DeltaNet
  gates, convolutions, norms, embedding and LM head in llama.cpp's own formats, and ships the decoder weights as new
  ggml types in a patched llama.cpp, so the models run in a standard server with tool calling, thinking and
  speculative decoding.
\end{itemize}
In one line: the ShiftAddLLM weight format, refitted with an output-aware, distillation-based pipeline, and deployed
through a LUT-free llama.cpp kernel.

\section{Why it beats the k-quant baselines}

\subsection{Memory}
Every decoder weight group is stored as $Q$ sign planes plus one u8 scale per plane and group and one fp16 per
256-weight super-block, i.e.\ about $Q + 0.3$ bits per weight for $G = 128$; with the per-linear plan the released
models average 4.4 and 4.0 bits per decoder weight. Against the reference Q4\_K\_M deployment (5.87\,GB) this saves
13\% (4.4 bits, 5.10\,GB) and 19\% (4.0 bits, 4.76\,GB) of file size and GPU-resident weight memory, so a 32K-context
server fits in 6--8\,GB of GPU memory. The 4.0-bit model is also smaller than Q3\_K\_M (4.83\,GB) while having
$1.7\times$ lower KLD.

\subsection{Fidelity: a better size--quality frontier}
Table~1 of \texttt{docs/results.md} gives, on the mixed-domain text, KLD $0.030$ for the 4.4-bit model (5.10\,GB)
against $0.036$ for UD-Q3\_K\_XL (5.24\,GB), and $0.038$ for the 4.0-bit model (4.76\,GB) against $0.067$ for
Q3\_K\_M (4.83\,GB): at every size from 3.9 to 5.1\,GB the ShiftAdd model has lower KLD than the k-quant of equal or
larger size. Four ingredients contribute:
\begin{enumerate}
\item \textbf{Output-aware fit} (\texttt{saq/gptq\_bcq.py}, \texttt{saq/colwise.py}). Each linear is quantized by a
  GPTQ column sweep with act-order against the layer Hessian; the scales of each group are fitted on the
  error-compensated weights and, after the sweep, all scales are re-solved exactly for the Hessian-weighted output
  error with the codes frozen. k-quants fit each block of 32 weights to a uniform grid by weight error alone.
\item \textbf{Mixed precision chosen by measurement} (\texttt{saq/sens\_kld.py}, \texttt{saq/sensitivity.py}). The
  end-to-end KLD increase of quantizing each linear alone at every (bits, group) setting is measured, and a knapsack
  selects the setting per linear for the target average, instead of one fixed recipe by tensor type for every layer.
\item \textbf{End-to-end scale distillation} (\texttt{saq/finetune.py}). With the codes frozen, every scale is trained as
  $\alpha_0(1+\theta)$ to match the BF16 teacher's top-64 log-probabilities on $768 \times 2048$ tokens, then snapped
  back to the u8 grid.
\item \textbf{A more expressive grid for the same bits.} A $Q$-plane group has $2^Q$ levels at $\sum_p \pm \alpha_p$,
  so their spacing is set by $Q$ independent scales and adapts to each group's weight distribution, with about 0.3 bits
  of scale overhead against 4.5 bits per weight for Q4\_K's uniform levels.
\end{enumerate}

\subsection{Speed}
\begin{itemize}
\item \textbf{Decode} (tg128 on the A100): 134.6 / 129.9 tok/s for the 4.0 / 4.4-bit models against 125.2 for
  Q4\_K\_M (+8\% / +4\%). Decode is memory-bound, so the smaller weights are read faster, and the mat-vec of
  Section~3 is multiplication-free: the weights enter only as dp4a byte masks, with one scale applied per group and
  plane, and no lookup tables.
\item \textbf{Prefill} (pp512): 5367 / 5525 tok/s against 4227 (+27\% / +31\%). The batched path dequantizes the planes
  to fp16 with Eq.~\eqref{eq:bcq} and runs the GEMM on the tensor cores through cuBLAS.
\item \textbf{MTP speculative decoding} keeps its acceptance rate (0.60--0.68) and reaches 171.5 / 183.3 tok/s with the
  draft head, against 158.2 for Q4\_K\_M.
\end{itemize}
\end{document}
